\section*{Chapitre \ref{ch:g12:proton-nmr} --- La RMN du proton}

\begin{solution}{exo:g12:proton-nmr:1}
Méthane~: 1. Éthane~: 1 (les deux \ce{CH3} sont identiques). Propane~: 2
(deux \ce{CH3} identiques, un \ce{CH2}). Méthanol~: 2 (\ce{CH3} et
\ce{OH}). Méthoxyméthane~: 1.
\end{solution}

\begin{solution}{exo:g12:proton-nmr:2}
2 voisins~: un triplet, 1:2:1. 3 voisins~: un quadruplet, 1:3:3:1.
0 voisin~: un singulet.
\end{solution}

\begin{solution}{exo:g12:proton-nmr:3}
Au total \qty{40}{mm} pour 8 \omterm{def:g9:inside-the-atom:proton}{protons}~: \qty{5}{mm} par \omterm{def:g9:inside-the-atom:proton}{proton}. Les
signaux représentent 3, 2 et 3 \omterm{def:g9:inside-the-atom:proton}{protons}.
\end{solution}

\begin{solution}{exo:g12:proton-nmr:4}
\omterm{def:g11:functional-groups:carbonyl}{Aldéhyde}~: de 9.7 à \qty{10.0}{ppm}. Un \ce{CH3} d'une chaîne~: de 0.7 à
\qty{1.3}{ppm}.
\end{solution}

\begin{solution}{exo:g12:proton-nmr:5}
Ses douze \omterm{def:g12:proton-nmr:equivalent}{protons équivalents} donnent une seule raie, fine et intense,
facile à placer à zéro. Le \omterm{def:g6:solutions-and-solubility:solute}{solvant} deutéré ne donne aucun signal de
\omterm{def:g9:inside-the-atom:proton}{proton} propre, qui sinon écraserait ceux des quelques milligrammes
d'échantillon.
\end{solution}

\begin{solution}{exo:g12:proton-nmr:6}
Deux signaux. Le \ce{CH2} voisin du chlore~: 2 \omterm{def:g9:inside-the-atom:proton}{protons}, un quadruplet (3
voisins), dans le domaine 2.5--\qty{4.0}{ppm}. Le \ce{CH3}~: 3 \omterm{def:g9:inside-the-atom:proton}{protons},
un triplet (2 voisins), près du domaine habituel des \ce{CH3}, un peu
plus haut parce que le chlore est à deux \omterm{def:g7:atoms-and-molecules:atom}{atomes} de distance.
\end{solution}

\begin{solution}{exo:g12:proton-nmr:7}
Trois signaux~: les deux \ce{CH3}, 6 \omterm{def:g12:proton-nmr:equivalent}{protons équivalents}, un doublet
(1 voisin, le \ce{CH})~; le \ce{CH}, 1 \omterm{def:g9:inside-the-atom:proton}{proton}, couplé à 6 voisins
donc 7 raies, dans le domaine \ce{H-C-O} (3.3--\qty{4.5}{ppm})~; le
\ce{OH}, 1 \omterm{def:g9:inside-the-atom:proton}{proton}, un singulet.
\end{solution}

\begin{solution}{exo:g12:proton-nmr:8}
(a) la propanone~; (b) l'éthanoate d'éthyle~; (c) l'acide propanoïque.
\end{solution}

\begin{solution}{exo:g12:proton-nmr:9}
Le \ce{CH2} est lié à l'\omterm{def:g7:atoms-and-molecules:atom}{atome} d'oxygène de l'\omterm{def:g11:functional-groups:carboxylic-acid}{ester}, électronégatif, qui
déblinde ses \omterm{def:g9:inside-the-atom:proton}{protons}~: domaine \ce{H-C-O},
3.3--\qty{4.5}{ppm}. Le \ce{CH3} est un carbone plus loin et reste dans
le domaine des chaînes.
\end{solution}

\begin{solution}{exo:g12:proton-nmr:10}
La propanone~: un \ce{C=O} (\omterm{def:g11:functional-groups:carbonyl}{cétone}, \qty{1715}{cm^{-1}}) et six \omterm{def:g12:proton-nmr:equivalent}{protons
équivalents}, un seul singulet. Le propanal donnerait trois signaux, dont
celui du \omterm{def:g9:inside-the-atom:proton}{proton} aldéhydique de \ce{CH3-CH2-CHO}, vers 9.7--\qty{10.0}{ppm}.
\end{solution}

\begin{solution}{exo:g12:proton-nmr:11}
Le \omterm{def:g9:inside-the-atom:proton}{proton} du \ce{OH} s'échange rapidement entre \omterm{def:g7:atoms-and-molecules:molecule}{molécules}~: il donne un
singulet et ne couple pas avec ses voisins. Le \ce{CH2} n'est donc couplé
qu'aux 3 \omterm{def:g9:inside-the-atom:proton}{protons} du \ce{CH3}~: un quadruplet.
\end{solution}

\begin{solution}{exo:g12:proton-nmr:12}
Acide butanoïque \ce{CH3-CH2-CH2-COOH}~: 4 signaux, un triplet
(\ce{CH3}), un signal à 6 raies (\ce{CH2} central, 5 voisins), un triplet
(\ce{CH2-C=O}), un singulet vers 11--\qty{12}{ppm}. Éthanoate d'éthyle~:
un singulet, un quadruplet, un triplet. Propanoate de méthyle
\ce{CH3-CH2-CO-O-CH3}~: un triplet, un quadruplet (\ce{CH2-C=O},
2.0--\qty{2.4}{ppm}), un singulet (\ce{O-CH3}). Méthanoate de propyle
\ce{H-CO-O-CH2-CH2-CH3}~: 4 signaux, un singulet très à gauche (le H
porté par le carbone du \ce{C=O}), un triplet (\ce{O-CH2}), un signal à
6 raies, un triplet (\ce{CH3}).
\end{solution}

\begin{solution}{exo:g12:proton-nmr:13}
Le signal du \ce{OH}~: les \omterm{def:g9:inside-the-atom:proton}{protons} \ce{O-H} s'échangent avec le
deutérium de l'eau lourde, et \ce{O-D} ne donne aucun signal de \omterm{def:g9:inside-the-atom:proton}{proton}.
Un signal qui disparaît après agitation avec \ce{D2O} appartient à un
\omterm{def:g9:inside-the-atom:proton}{proton} porté par O ou N.
\end{solution}

\begin{solution}{exo:g12:proton-nmr:14}
L'éthanol ajoute un quadruplet vers \qty{3.69}{ppm}, un triplet vers
\qty{1.23}{ppm} et un singulet du \ce{OH}. Le singulet de l'\omterm{def:g11:functional-groups:carboxylic-acid}{ester}~:
\qty{30}{mm} pour 3 \omterm{def:g9:inside-the-atom:proton}{protons}, soit \qty{10}{mm} par \omterm{def:g9:inside-the-atom:proton}{proton} d'\omterm{def:g11:functional-groups:carboxylic-acid}{ester}. Le
triplet de l'éthanol~: \qty{3}{mm} pour 3 \omterm{def:g9:inside-the-atom:proton}{protons}, soit \qty{1}{mm} par
\omterm{def:g9:inside-the-atom:proton}{proton} d'éthanol. Comme chaque \omterm{def:g7:atoms-and-molecules:molecule}{molécule} apporte un \ce{CH3} à ces
signaux, les quantités sont dans le rapport $1/10$~: environ une
\omterm{def:g7:atoms-and-molecules:molecule}{molécule} d'éthanol pour dix d'\omterm{def:g11:functional-groups:carboxylic-acid}{ester}.
\end{solution}

\begin{solution}{exo:g12:proton-nmr:15}
Quatre signaux~: les deux \ce{CH3} (6 \omterm{def:g9:inside-the-atom:proton}{protons}), un doublet~; le \ce{CH}
(1 \omterm{def:g9:inside-the-atom:proton}{proton}), couplé à 6 + 2 = 8 voisins, donc 9 raies~; le \ce{CH2} (2
\omterm{def:g9:inside-the-atom:proton}{protons}) voisin de \ce{O}, un doublet, dans le domaine \ce{H-C-O}~; le
\ce{OH} (1 \omterm{def:g9:inside-the-atom:proton}{proton}), un singulet. C'est le signal du \ce{CH} qui compte
le plus de raies.
\end{solution}

\begin{solution}{pb:g12:proton-nmr:1}
\textbf{1.} Un \omterm{def:g11:organic-skeletons:alkane}{alcane} à quatre carbones a pour formule \ce{C4H10}~; un
\ce{C=O} retire deux hydrogènes et les deux oxygènes n'en ajoutent
aucun~: \ce{C4H8O2}, la formule des acides et des \omterm{def:g11:functional-groups:carboxylic-acid}{esters} à quatre
carbones.

\textbf{2.} \ce{CH3-CH2-CH2-COOH}, acide butanoïque~;
\ce{CH3-CO-O-CH2-CH3}, éthanoate d'éthyle~; \ce{CH3-CH2-CO-O-CH3},
propanoate de méthyle~; \ce{H-CO-O-CH2-CH2-CH3}, méthanoate de propyle.

\textbf{3.} Les trois \omterm{def:g11:functional-groups:carboxylic-acid}{esters} ont en commun le groupe \omterm{def:g11:functional-groups:carboxylic-acid}{ester}~; les quatre
ont un \ce{C=O}.

\textbf{4.} Un \ce{C=O}, dans le domaine des \omterm{def:g11:functional-groups:carboxylic-acid}{esters} (vers
\qty{1735}{cm^{-1}}).

\textbf{5.} La très large bande \ce{O-H} de 2500 à
\qty{3300}{cm^{-1}}~: elle est absente, ce n'est donc pas l'acide
butanoïque.

\textbf{6.} Non~: les trois \omterm{def:g11:functional-groups:carboxylic-acid}{esters} ont les mêmes groupes.

\textbf{7.} Trois.

\textbf{8.} Un groupe éthyle \ce{CH3-CH2-}~: le \ce{CH2} couplé à 3
\omterm{def:g9:inside-the-atom:proton}{protons} (quadruplet), le \ce{CH3} à 2 (triplet).

\textbf{9.} Ses \omterm{def:g9:inside-the-atom:proton}{protons} n'ont aucun \omterm{def:g9:inside-the-atom:proton}{proton} sur les \omterm{def:g7:atoms-and-molecules:atom}{atomes} voisins~: c'est
un \ce{CH3} lié au carbone du \ce{C=O} ou à l'oxygène du cycle.

\textbf{10.} Le \ce{CH2} est lié à l'oxygène (domaine \ce{H-C-O}).

\textbf{11.} Un triplet (\ce{CH3}), un quadruplet vers 2.0--\qty{2.4}{ppm}
(\ce{CH2} voisin du \ce{C=O}) et un singulet dans le domaine \ce{H-C-O}
(\ce{O-CH3}). Le quadruplet serait bien en dessous de \qty{4.12}{ppm} et
le singulet bien au-dessus de \qty{2.04}{ppm}~: ce n'est pas l'inconnu.

\textbf{12.} Quatre signaux (un singulet très à gauche, un triplet, un
signal à 6 raies, un triplet)~: ce n'est pas l'inconnu, qui en a trois.

\textbf{13.} L'éthanoate d'éthyle, \ce{CH3-CO-O-CH2-CH3}~: singulet à
\qty{2.04}{ppm} = \ce{CH3-C=O}~; quadruplet à \qty{4.12}{ppm} =
\ce{O-CH2}~; triplet à \qty{1.26}{ppm} = le \ce{CH3} terminal.

\textbf{14.} Singulet 3, quadruplet 2, triplet 3.

\textbf{15.} $72 / 8 = \qty{9}{mm}$ par \omterm{def:g9:inside-the-atom:proton}{proton}.

\textbf{16.} Singulet~: $3 \times 9 = \qty{27}{mm}$~; triplet~:
\qty{27}{mm}.

\textbf{17.} $18 + 27 + 27 = \qty{72}{mm}$.

\textbf{18.} Oui~: les petits \omterm{def:g11:functional-groups:carboxylic-acid}{esters} sont connus pour leurs odeurs
fruitées.

\textbf{19.} La marche du quadruplet mesure \qty{18}{mm} sur
\qty{72}{mm}.
\end{solution}
